% ============================================================================
% WORKBOOK — Additional: Introduction to Scientific Notation
% State-specific (VA)
% Practice-focused reinforcement for the study guide topic
% ============================================================================
\section{Introduction to Scientific Notation}
\topicTitle{Introduction to Scientific Notation}

\begin{quickReview}{Scientific Notation}
Scientific notation writes a number as $a \times 10^n$ where:
\begin{itemize}[leftmargin=*]
    \item $a$ is at least $1$ but less than $10$ (one non-zero digit before the decimal)
    \item $n$ is an integer (positive for large numbers, negative for small numbers)
    \item \textbf{Large numbers:} move the decimal \emph{left} to get $a$; $n$ is positive.
    \item \textbf{Small numbers:} move the decimal \emph{right} to get $a$; $n$ is negative.
\end{itemize}

\medskip
\textbf{Examples:} $45{,}000 = 4.5 \times 10^4$ \qquad $0.003 = 3 \times 10^{-3}$
\end{quickReview}

% ── Warm-Up ──────────────────────────────────────────────────────────────────
\begin{practiceBox}[funGreen]{\faSun~Warm-Up}

\practiceHeader[funGreen]{Powers of 10}

\begin{multicols}{2}
\prob $10^3 =$ \answerBlank[2cm]
\answerExplain{$1{,}000$}{$10^3 = 10 \times 10 \times 10 = 1{,}000$.}

\prob $10^5 =$ \answerBlank[2cm]
\answerExplain{$100{,}000$}{$10^5$ means a $1$ followed by $5$ zeros: $100{,}000$.}

\prob $10^1 =$ \answerBlank[2cm]
\answerExplain{$10$}{$10^1 = 10$.}

\prob $10^0 =$ \answerBlank[2cm]
\answerExplain{$1$}{Any nonzero number to the zero power equals $1$.}

\prob $10^{-2} =$ \answerBlank[2cm]
\answerExplain{$0.01$}{$10^{-2} = \frac{1}{10^2} = \frac{1}{100} = 0.01$.}

\prob $10^{-4} =$ \answerBlank[2cm]
\answerExplain{$0.0001$}{$10^{-4} = \frac{1}{10^4} = \frac{1}{10{,}000} = 0.0001$.}
\end{multicols}
\end{practiceBox}

% ── Practice A: Standard to Scientific ───────────────────────────────────────
\begin{practiceBox}{Convert to Scientific Notation}

\begin{multicols}{2}
\prob $93{,}000{,}000$ \answerBlank[3cm]
\answerExplain{$9.3 \times 10^7$}{Move the decimal $7$ places left to get $9.3$. So $93{,}000{,}000 = 9.3 \times 10^7$.}

\prob $5{,}400$ \answerBlank[3cm]
\answerExplain{$5.4 \times 10^3$}{Move the decimal $3$ places left: $5.4 \times 10^3$.}

\prob $720{,}000$ \answerBlank[3cm]
\answerExplain{$7.2 \times 10^5$}{Move the decimal $5$ places left: $7.2 \times 10^5$.}

\prob $0.00056$ \answerBlank[3cm]
\answerExplain{$5.6 \times 10^{-4}$}{Move the decimal $4$ places right to get $5.6$. The exponent is $-4$ because the number is small.}

\prob $0.0082$ \answerBlank[3cm]
\answerExplain{$8.2 \times 10^{-3}$}{Move the decimal $3$ places right: $8.2 \times 10^{-3}$.}

\prob $0.0000039$ \answerBlank[3cm]
\answerExplain{$3.9 \times 10^{-6}$}{Move the decimal $6$ places right: $3.9 \times 10^{-6}$.}
\end{multicols}
\end{practiceBox}

% ── Practice B: Scientific to Standard ───────────────────────────────────────
\begin{practiceBox}[funTeal]{Convert to Standard Form}

\begin{multicols}{2}
\prob $3.07 \times 10^5$ \answerBlank[3cm]
\answerExplain{$307{,}000$}{Move the decimal $5$ places right: $3.07 \to 307{,}000$.}

\prob $6.5 \times 10^3$ \answerBlank[3cm]
\answerExplain{$6{,}500$}{Move the decimal $3$ places right: $6.5 \to 6{,}500$.}

\prob $1.2 \times 10^8$ \answerBlank[3cm]
\answerExplain{$120{,}000{,}000$}{Move the decimal $8$ places right: $120{,}000{,}000$.}

\prob $4.1 \times 10^{-2}$ \answerBlank[3cm]
\answerExplain{$0.041$}{Move the decimal $2$ places left: $4.1 \to 0.041$.}

\prob $7.03 \times 10^{-5}$ \answerBlank[3cm]
\answerExplain{$0.0000703$}{Move the decimal $5$ places left: $7.03 \to 0.0000703$.}

\prob $9.9 \times 10^6$ \answerBlank[3cm]
\answerExplain{$9{,}900{,}000$}{Move the decimal $6$ places right: $9.9 \to 9{,}900{,}000$.}
\end{multicols}
\end{practiceBox}

% ── Practice C: Comparing ────────────────────────────────────────────────────
\begin{practiceBox}[funOrange]{Which Is Larger?}

\practiceHeader[funOrange]{Write $>$, $<$, or $=$ between the two numbers.}

\prob $6.1 \times 10^8$ \answerBlank[1cm] $9.2 \times 10^7$
\answerExplain{$>$}{Compare exponents first: $10^8 > 10^7$. Since $6.1 \times 10^8 = 610{,}000{,}000$ and $9.2 \times 10^7 = 92{,}000{,}000$, the first is larger.}

\prob $3.5 \times 10^4$ \answerBlank[1cm] $8.1 \times 10^4$
\answerExplain{$<$}{Same exponent, so compare the first numbers: $3.5 < 8.1$.}

\prob $2.0 \times 10^{-3}$ \answerBlank[1cm] $5.0 \times 10^{-4}$
\answerExplain{$>$}{$10^{-3} > 10^{-4}$. So $2.0 \times 10^{-3} = 0.002$ and $5.0 \times 10^{-4} = 0.0005$. The first is larger.}

\prob $7.7 \times 10^5$ \answerBlank[1cm] $7.7 \times 10^6$
\answerExplain{$<$}{Same coefficient but $10^5 < 10^6$, so $7.7 \times 10^5 < 7.7 \times 10^6$.}
\end{practiceBox}

% ── True or False ────────────────────────────────────────────────────────────
\begin{practiceBox}[funPurple]{True or False?}

\trueOrFalse{$45 \times 10^3$ is written in correct scientific notation.}
\answerExplain{False}{In scientific notation, $a$ must be at least $1$ and less than $10$. Since $45 \geq 10$, this is not correct. It should be $4.5 \times 10^4$.}

\trueOrFalse{A negative exponent means the number is less than $1$.}
\answerExplain{True}{A negative exponent in $a \times 10^n$ means the number is a small decimal less than $1$.}
\end{practiceBox}

% ── Word Problems ────────────────────────────────────────────────────────────
\begin{practiceBox}[funBlue]{\faGlobe~Word Problems}

\wordProblem{The Sun is about $1.5 \times 10^8$ km from Earth. Write this in standard form.}{km}
\answerExplain{$150{,}000{,}000$ km}{Move the decimal $8$ places right: $1.5 \to 150{,}000{,}000$.}

\wordProblem{A red blood cell is about $0.000007$ m in diameter. Write this in scientific notation.}{m}
\answerExplain{$7 \times 10^{-6}$ m}{Move the decimal $6$ places right to get $7$. The exponent is $-6$.}
\end{practiceBox}

% ── Challenge ────────────────────────────────────────────────────────────────
\begin{challengeBox}
\prob The distance from Earth to the nearest star (Proxima Centauri) is about $4.01 \times 10^{13}$ km. The distance from Earth to the Moon is about $3.84 \times 10^5$ km. About how many times farther is Proxima Centauri than the Moon? \answerBlank[3cm]
\answerExplain{About $1.04 \times 10^8$ times farther}{Divide: $\frac{4.01 \times 10^{13}}{3.84 \times 10^5} \approx 1.04 \times 10^{13-5} = 1.04 \times 10^8$, or about $104$ million times farther.}
\end{challengeBox}

\encouragement{Fantastic work with scientific notation! You can now express any number --- from the tiniest atom to the biggest galaxy!}
